\chapter{Gestational Trophoblastic Disease}
\index{Gestational Trophoblastic Disease}
\label{sec:Gestational Trophoblastic Disease}

\section{Overview}
Gestational trophoblastic disease (GTD) encompasses a spectrum of disorders arising from abnormal trophoblast growth, ranging from non-cancerous hydatidiform mole to cancerous choriocarcinoma. Incidence is 1 in 1000 pregnancies in Western countries. This pregnancy-related condition requires careful monitoring to ensure the safety of both mother and baby. Early detection and management are essential.
\section{Causes}
\section{Risk Factors}

\begin{itemize}[noitemsep]
  \item Genetic predisposition and family history
  \item Advancing age
  \item Lifestyle factors (diet, exercise, smoking, alcohol)
  \item Environmental exposures
  \item Underlying medical conditions
  \item Medication use
\end{itemize}
\section{Signs and Symptoms}

\subsection{Common symptoms}
\begin{itemize}[noitemsep]
  \item You may experience vaginal bleeding, exaggerated uterine size for dates, hyperemesis gravidarum, and elevated $\beta$-hCG. Ultrasonography shows a characteristic ``snowstorm'' pattern in complete mole. 
  \end{itemize}

\subsection{Emergency warning signs}
\begin{itemize}[noitemsep]
  \item Severe or worsening symptoms of gestational trophoblastic disease require immediate medical evaluation
\end{itemize}
\section{Progression and Stages}
Hydatidiform mole is classified as complete (46,XX, entirely paternal genome, risk of cancerous transformation 15--20\%) or partial (69,XXY, triploid, risk 1--5\%). The condition may progress through different stages of severity over time. Early stages often involve mild or subtle symptoms that may be overlooked. Moderate stages involve more noticeable symptoms affecting daily function. Advanced stages may involve significant impairment and complications.
\section{Complications}
\begin{itemize}[noitemsep]
  \item Disease progression leading to worsening symptoms
  \item Secondary infections or injuries
  \item Side effects of treatment
  \item Reduced quality of life
  \item Emotional and psychological impact
\end{itemize}
\section{Diagnosis}
Comprehensive medical history and physical examination Laboratory tests to assess organ function and rule out other conditions Imaging studies to visualize affected structures Specialized testing depending on the affected system Consultation with relevant specialists Monitoring of disease progression over time

\begin{itemize}[noitemsep]
  \item Comprehensive medical history and physical examination
  \item Laboratory tests to assess organ function
  \item Imaging studies to visualize affected structures
  \item Specialized testing depending on the affected system
  \item Consultation with relevant specialists
\end{itemize}
\section{Prevention}
\begin{itemize}[noitemsep]
  \item Maintain a healthy lifestyle with balanced diet and exercise
  \item Avoid known risk factors
  \item Attend regular medical check-ups for early detection
  \item Follow recommended screening guidelines
\end{itemize}
\section{Treatment Options}
Treatment typically involves suction dilation and curettage with serial $\beta$-hCG monitoring. Gestational trophoblastic neoplasia (GTN) is diagnosed when $\beta$-hCG plateaus or rises after molar evacuation

\begin{itemize}[noitemsep]
  \item Low-risk GTN doctors typically treat it with single-agent methotrexate or actinomycin D
  \item High-risk GTN requires multi-agent chemotherapy.
  \item Medications to manage symptoms and address underlying mechanisms
  \item Lifestyle modifications and self-care measures
  \item Physical therapy and rehabilitation
  \item Surgical intervention when indicated
  \item Regular monitoring and follow-up care
\end{itemize}
\section{Living with It}
\begin{itemize}[noitemsep]
  \item Follow your treatment plan as prescribed
  \item Maintain regular follow-up with your healthcare provider
  \item Make necessary lifestyle adjustments
  \item Seek support from family, friends, or support groups
  \item Stay informed about your condition
\end{itemize}
\section{Outlook}
Most people recover well, with cure rates exceeding 90\% for low-risk and 80--85\% for high-risk disease. The outlook varies depending on the specific condition, severity at diagnosis, and response to treatment. Early intervention generally leads to better outcomes. Many conditions can be effectively managed with appropriate treatment. Regular follow-up care is important for monitoring and treatment adjustment.
